\chapter{Diseño experimental e implementación}\label{chap:experimentos}

% Barrera de flotantes: evita que las tablas migren de sección. Si el paquete
% placeins esta cargado se usa su \FloatBarrier; si no, se recurre a \clearpage.
\providecommand{\FloatBarrier}{\clearpage}

Este capítulo concreta la batería de experimentos ejecutada sobre los tres \textit{frameworks}
con la metodología descrita en el Capítulo~\ref{chap:metodologia}. Se define primero el escenario
de entrenamiento común a las tres herramientas —conjunto de datos, modelo y configuración de
aprendizaje—, que se mantiene constante para que las diferencias observadas se atribuyan al
\textit{framework} y no a otros factores. A continuación se presentan la matriz de experimentos y
los dos modos de ejecución, secuencial y concurrente, en los que se ejecuta cada configuración.

\section{Diseño del experimento común}\label{sec:diseno-comun}

La comparación parte de un principio básico de control experimental: todo aquello que no es el
propio \textit{framework} debe permanecer constante. Por ese motivo, el conjunto de datos, el
modelo, el algoritmo de aprendizaje federado, los hiperparámetros y el
hardware son idénticos en las tres herramientas, y el particionamiento de los datos es
equivalente entre ellas: misma naturaleza IID y mismos tamaños por cliente, aunque cada
\textit{framework} lo materialice con su propia tubería de carga de datos
(Sección~\ref{sec:datos}). La única variable que cambia entre ejecuciones
es la herramienta evaluada y, dentro de cada una, el modo de ejecución descrito en la
Sección~\ref{sec:modos-ejecucion}.

La implementación del experimento en cada \textit{framework} se apoyó en su documentación
oficial y en los proyectos de ejemplo que proporciona cada herramienta, tomados
como punto de partida y adaptados al escenario común descrito en este capítulo.
En Flower se siguieron su guía y ejemplos oficiales \parencite{beutel2020flower}; en
NVIDIA FLARE, la documentación del SDK y sus tutoriales \parencite{roth2022nvflare}; y
en FedML, la documentación de la plataforma y sus implementaciones de referencia
\parencite{he2020fedml}.

La implementación se reparte de forma deliberada en dos planos. Este capítulo recoge el
plano común y operativo —el escenario idéntico para las tres herramientas y las órdenes
concretas con que se ejecuta cada modo (Sección~\ref{sec:modos-ejecucion})—, mientras que
las adaptaciones de código propias de cada \textit{framework} —modelo, carga de datos,
bucle de entrenamiento, estrategia de agregación y automatización de la ejecución— se
documentan junto a la evaluación de flexibilidad de cada herramienta, en las
Secciones~\ref{sec:flex-flower}, \ref{sec:flex-nvflare} y~\ref{sec:flex-fedml}. Se ha
optado por esta organización para mantener cada implementación al lado del análisis que la
interpreta y evitar duplicar el código entre capítulos.

\subsection{Conjunto de datos y particionamiento}\label{sec:datos}

Como banco de pruebas se emplea CIFAR-10, un conjunto de 60\,000 imágenes en color de
$32\times32$ píxeles distribuidas en diez clases. Se respeta la partición oficial de
50\,000 imágenes de entrenamiento y 10\,000 de prueba. Las 50\,000 imágenes de entrenamiento se
dividen, a su vez, en 40\,000 para el entrenamiento federado y 10\,000 para validación local en
cada cliente, mientras que las 10\,000 imágenes de prueba se reservan para la evaluación
centralizada del modelo global, ajena al entrenamiento.

El reparto entre clientes es independiente e idénticamente distribuido (IID). Para ello se
barajan los índices con una semilla fija y se asignan en bloques de igual tamaño a cada cliente,
de modo que todos reciben una muestra estadísticamente equivalente del conjunto global. Conviene
matizar que cada \textit{framework} construye ese reparto con su propia tubería de carga de
datos —Flower mediante el particionador IID de \texttt{flwr-datasets} y NVIDIA FLARE y FedML
mediante un barajado global con \texttt{torchvision}—, por lo que las particiones no contienen
necesariamente las mismas imágenes en cada cliente, pero sí comparten la naturaleza IID y los
mismos tamaños por cliente; las diferencias observadas no son, por tanto, atribuibles al reparto.
Se ha
optado por una distribución IID de forma deliberada: el objetivo es aislar el comportamiento del
\textit{framework} del efecto añadido de la heterogeneidad de datos, que introduciría una
fuente de variación difícil de separar. El estudio de escenarios no~IID queda planteado como
ampliación futura y se recoge en la Sección~\ref{sec:alcance-limitaciones}.

Sobre las imágenes se aplican las transformaciones habituales para CIFAR-10. En entrenamiento se
usa aumento de datos mediante recorte aleatorio con relleno de cuatro píxeles y volteo
horizontal aleatorio; en validación y prueba no se aplica aumento. En todos los casos las
imágenes se normalizan con las medias y desviaciones típicas por canal propias del conjunto.

\subsection{Modelo}\label{sec:modelo}

El modelo común es una red residual ResNet-18 adaptada a las dimensiones de CIFAR-10. La
arquitectura original está pensada para imágenes de $224\times224$ píxeles, por lo que se
sustituye la primera convolución por una de núcleo $3\times3$ con paso 1 y se elimina la capa de
\textit{maxpooling} inicial, de manera que la red no reduce en exceso la resolución de unas
imágenes ya pequeñas. La capa totalmente conectada final se ajusta a diez salidas, una por
clase. Se eligió ResNet-18 por tratarse de una arquitectura estándar, suficientemente expresiva
para esta tarea y con un coste de cómputo asumible en una GPU de 6\,GB. El modelo se implementa
de forma idéntica en los tres \textit{frameworks}.

\subsection{Algoritmo y configuración de entrenamiento}\label{sec:config-entrenamiento}

El algoritmo de aprendizaje federado es Federated Averaging (FedAvg), descrito en el capítulo de
fundamentos teóricos. En cada ronda, los clientes entrenan el modelo global con sus datos
locales y el servidor promedia los modelos locales resultantes, ponderándolos por el número de
ejemplos de cada cliente. El entrenamiento local emplea descenso de gradiente estocástico (SGD) con tasa de
aprendizaje 0{,}02, momento 0{,}9 y \textit{weight decay} de $5\times10^{-4}$, con un tamaño de
lote de 32 y entropía cruzada como función de pérdida. Estos valores son comunes a las tres
herramientas.

Los hiperparámetros no se determinaron mediante una búsqueda exhaustiva, sino que se fijaron a
priori para disponer de una configuración común, controlada y comparable entre los tres
\textit{frameworks}. El objetivo del trabajo no es maximizar el rendimiento del modelo, sino
comparar Flower, NVIDIA FLARE y FedML bajo idénticas condiciones; por ese motivo se escogieron
valores moderados y estables, atendiendo a configuraciones habituales de FedAvg, al hardware
disponible y a una serie de pruebas preliminares que confirmaron que los experimentos se
ejecutaban correctamente. Esta decisión se asume como una limitación del estudio
(Sección~\ref{sec:alcance-limitaciones}): los resultados corresponden a esta configuración
concreta y no a una optimización individual de cada herramienta.

\section{Matriz de experimentos}\label{sec:matriz-experimentos}

A partir del escenario común se define una matriz de experimentos que varía de forma controlada
el número de clientes, el número de rondas, las épocas locales y la fracción de clientes que
participa en cada ronda. La Tabla~\ref{tab:matriz-experimentos} resume las configuraciones. El
experimento e0 actúa como referencia; e1 constituye la configuración principal del estudio, con
más épocas locales por ronda; e2 explora la escalabilidad variando el número de clientes; y e3
analiza la participación parcial, en la que solo una fracción de los clientes interviene en cada
ronda.

\begin{table}[htbp]
  \centering
  \small
  \caption{Matriz de experimentos del estudio. Cada configuración se ejecuta en los tres
  \textit{frameworks} y en los dos modos de ejecución.}
  \label{tab:matriz-experimentos}
  \begin{tabular}{llccccc}
    \toprule
    ID & Propósito & Clientes & Clientes/ronda & Rondas & Épocas locales \\
    \midrule
    e0 & Línea base de referencia      & 10 & 10 & 30 & 2  \\
    e1 & Configuración principal       & 5  & 5  & 20 & 10 \\
    e2-c4 & Escalabilidad (4 clientes)    & 4  & 4  & 10 & 2  \\
    e2-c5 & Escalabilidad (5 clientes)    & 5  & 5  & 10 & 2  \\
    e2-c8 & Escalabilidad (8 clientes)    & 8  & 8  & 10 & 2  \\
    e3 & Participación parcial         & 8  & 4  & 20 & 10 \\
    \bottomrule
  \end{tabular}
\end{table}

El experimento e3 exige una aclaración terminológica. Mide la \textit{participación parcial}
de clientes, es decir, el efecto de que solo una fracción del total intervenga en cada ronda. No
debe confundirse con la \textit{flexibilidad} de cada herramienta, entendida como propiedad
arquitectónica, que se analiza de forma cualitativa en la Sección~\ref{sec:criterios-flexibilidad}.
A lo largo de la memoria, el término flexibilidad se reserva para ese eje cualitativo.

\section{Modos de ejecución: secuencial y concurrente}\label{sec:modos-ejecucion}

Cada configuración de la matriz se ejecuta en dos modos distintos. En el modo \textit{secuencial}
se mantiene una sola tarea de entrenamiento activa cada vez, de manera que los clientes de una
ronda se procesan uno tras otro. En el modo \textit{concurrente} se intenta mantener varios
clientes o procesos activos a la vez, con el fin de observar cómo gestiona cada \textit{framework}
el paralelismo y qué impacto tiene en el consumo de recursos y en la estabilidad de la ejecución.

La forma de activar cada modo depende de la arquitectura de la herramienta. En NVIDIA FLARE, el
simulador admite un parámetro que fija el número de hilos de entrenamiento activos: con un único
hilo la ejecución es secuencial y, al aumentarlo hasta el número de clientes por ronda, se vuelve
concurrente. El Código~\ref{cod:nvflare-modo} muestra ambas invocaciones.

\begin{bashcodeNN}[title={Modos de ejecución en NVIDIA FLARE}, label={cod:nvflare-modo}, list entry={\protect\numberline{\thelisting}Modos de ejecución en NVIDIA FLARE}]
# Modo secuencial: un único hilo de entrenamiento activo
nvflare simulator -w workspace -n 5 -t 1 --gpu 0 job_dir

# Modo concurrente: tantos hilos como clientes por ronda
nvflare simulator -w workspace -n 5 -t 5 --gpu 0 job_dir
\end{bashcodeNN}

En FedML, el modo secuencial corresponde al backend de proceso único (\texttt{sp}), que entrena a
los clientes de forma sucesiva, mientras que el modo concurrente se apoya en el backend MPI,
lanzado con \texttt{mpirun}, que ejecuta varios procesos de cliente en paralelo
(Código~\ref{cod:fedml-modo}).

\begin{bashcodeNN}[title={Modos de ejecución en FedML}, label={cod:fedml-modo}, list entry={\protect\numberline{\thelisting}Modos de ejecución en FedML}]
# Modo secuencial: backend de proceso único
python main.py --cf config_sp.yaml

# Modo concurrente: backend MPI con varios procesos de cliente
mpirun --oversubscribe -np 6 python main.py --cf config_mpi.yaml
\end{bashcodeNN}

En Flower, la simulación se gestiona internamente mediante Ray y la ejecución se lanza con
\texttt{flwr run}, indicando la configuración de la federación. A diferencia de NVIDIA FLARE y
FedML, la concurrencia no se controla con un parámetro de hilos o de \textit{backend}, sino con
los recursos de GPU que se reservan para cada cliente mediante
\texttt{client-resources-num-gpus}. Si cada cliente reserva la GPU completa ($1{,}0$), solo puede
entrenar un cliente a la vez y la ejecución es secuencial; si reserva una fracción ($0{,}5$),
varios clientes caben simultáneamente en la GPU y la ejecución se vuelve concurrente. El número
total de clientes simulados se fija con \texttt{num-supernodes}. El Código~\ref{cod:flower-modo}
muestra ambas invocaciones.

\begin{bashcodeNN}[title={Modos de ejecución en Flower}, label={cod:flower-modo}, list entry={\protect\numberline{\thelisting}Modos de ejecución en Flower}]
# Modo secuencial: cada cliente reserva la GPU completa,
# por lo que solo se entrena un cliente a la vez
flwr run . --federation-config="num-supernodes=5 client-resources-num-gpus=1.0"

# Modo concurrente: cada cliente reserva una fraccion de la GPU,
# permitiendo varios clientes simultaneos
flwr run . --federation-config="num-supernodes=5 client-resources-num-gpus=0.5"
\end{bashcodeNN}
